% Generated draft: requires mathematical and source review.
\documentclass[11pt]{article}
\usepackage[margin=25mm]{geometry}
\usepackage{fontspec}
\setmainfont{Latin Modern Roman}
\usepackage{amsmath,amssymb,mathtools,unicode-math}
\setmathfont{Latin Modern Math}
\usepackage{enumitem,needspace,xurl,hyperref}
\setlength{\parindent}{0pt}
\setlength{\parskip}{5pt}
\setlength{\emergencystretch}{3em}
\newcommand{\N}{\mathbb{N}}
\newcommand{\Z}{\mathbb{Z}}
\newcommand{\Q}{\mathbb{Q}}
\newcommand{\R}{\mathbb{R}}
\newcommand{\C}{\mathbb{C}}
\newcommand{\F}{\mathbb{F}}
\newcommand{\A}{\mathbb{A}}
\newcommand{\PP}{\mathbb{P}}
\newcommand{\E}{\mathbb{E}}
\newcommand{\eps}{\varepsilon}
\newcommand{\dd}{\,\mathrm{d}}
\begin{document}
\section*{UnsolvedMath ID: 10700001}
\textbf{Source problem code:} catalogue entry
\subsection{Definitions and mathematical statement}
For a manifold \(M\), its free loop space is \(LM=\operatorname{Map}(S^1,M)\) with an appropriate mapping-space topology. In the usual closed oriented \(4\)-manifold setting, string topology includes the loop product
\[
 H_p(LM;R)\otimes H_q(LM;R)
       \longrightarrow H_{p+q-4}(LM;R),
\]
and further operations such as those induced by rotating loops. After the degree shift \(\mathbb H_k(LM;R)=H_{k+4}(LM;R)\), the product has degree zero. The coefficient ring and the collection of operations are not specified by the source.

Oleg Viro's Section 9 problem asks whether string topology is sensitive to smooth structures on \(4\)-manifolds. A usual existence interpretation asks whether two homeomorphic \(4\)-manifolds with distinct smooth structures can be distinguished by the intended string-topology data. A precise version must say which operations and extra structures are retained and whether comparison means abstract algebraic isomorphism or compatibility with a chosen underlying homeomorphism. The record gives no compactness, boundary or orientability assumptions, so the closed oriented loop-product example does not restrict the whole question without source review.

The supplied research summary reports a negative answer for a homotopy-level knot/embedding-space variant, attributing it to Knudsen--Kupers (2024). It separately expects insensitivity for the literal string-topology operation-level question, citing homotopy invariance of free-loop-space string topology for closed \(4\)-manifolds, but explicitly says the exact statement was not independently verified in full. The homotopy-level report is therefore not silently treated as a verified solution of the operation-level question.
\subsection{Short English statement}
Can string-topology operations distinguish different smooth structures on the same topological four-manifold? The relevant operations, manifold hypotheses and notion of comparison are unspecified. The supplied partial-progress assessment distinguishes a reported negative embedding-space result from an expected but not fully verified operation-level answer.
\subsection{Sources}
\begin{itemize}
\item Local UnsolvedMath registry, permanent ID \(10700001\), exact source code \texttt{catalogue entry}, ``String topology and smooth structures on 4-manifolds,'' attributed to Oleg Viro; no proposed year supplied.
\item Viro, \emph{Space of smooth 1-knots in a 4-manifold}, explicitly labeled problem in Section 9: \url{https://amj.math.stonybrook.edu/html-articles/Files-2000-2014/14-07/}. The article has one explicitly labeled problem in its string-topology section.
\item Supplied research summary names Knudsen--Kupers (2024) for homotopy-level knot/embedding-space insensitivity, without a complete citation. Its operation-level expectation is expressly not independently verified in full, and no separate primary citation for the stated string-topology homotopy invariance is supplied.
\item HTML-text extraction, accessed 2026-07-29. Source evidence NEEDS\_REVIEW presents an open problem and calls for a current primary-literature audit. Registry status partially\_solved and classification PARTIAL-PROGRESS retain the qualified assessment above.
\item Dataset attribution: Ulam AI / UnsolvedMath Contributors, CC BY 4.0. Source rights NEEDS\_REVIEW: publicly accessible author/publisher copy, redistribution terms requiring review.
\end{itemize}
\end{document}
